% =====================================================================
\chapter{Référence du contrat canonique}
\label{ann:contrat}
% =====================================================================

Cette annexe reprend la spécification complète du contrat de données, telle
qu'elle figure dans la spécification du projet. Elle est destinée à un
intégrateur qui souhaite consommer l'API.

\section{Structure complète}

\begin{lstlisting}[style=json,caption={Contrat canonique \texttt{DiagnosticRecord}
v1.0.},label={lst:contrat-full}]
POST /diagnostic/adresse
body: { "adresse": "10 rue de Rivoli, 75004 Paris" }

200 ->
{
  "schema_version": "1.0",
  "adresse": {
    "saisie": "...",              // entree brute
    "normalisee": "...",          // libelle du geocodeur
    "citycode": "75056",
    "postcode": "75004",
    "city": "Paris",
    "lat": 48.8566,
    "lon": 2.3522,
    "geocode_score": 0.95
  },
  "aleas": [                      // 13 objets, referentiel complet
    {
      "code": "inondation",
      "libelle": "Inondation",
      "present": true,            // true | false | null
      "present_commune": true,
      "zonage": "Dans un perimetre PPR inondation",   // texte verbatim
      "hauteur_eau_m": 1.5,       // optionnels specifiques a l'alea,
      "zone_sismique": null,      // TOUJOURS presents, null si non applicable
      "catnat_historique": [],
      "per_building": {           // present uniquement si teste geometriquement
        "method": "point-in-polygon",  // "point-in-polygon" | "proximity"
        "radius_m": null,              // echo de la proximite si applicable
        "count": 2
      },
      "source": "georisques",
      "url_detail": "https://www.georisques.gouv.fr/risques/inondations",
      "erreur": null,             // string | null
      "resolution": "per-building"
      // "per-building" | "commune-level" | "commune-level-estimate"
    }
  ],
  "bdnb": {                       // bloc entier null si BDNB indisponible
    "donnees": {
      // ~139 champs BDNB verbatim
      // batiment = batiment principal
      // autres_batiments_meme_adresse = [...]
    },
    "_source": {
      "provider": "BDNB",
      "url": "https://api.bdnb.io",
      "attribution": "Source BDNB (CSTB) - donnees a jour au <millesime>",
      "recuperee_le": "2026-08-25T10:00:00+00:00"
    }
  },
  "georisques_source": {
    "provider": "Georisques",
    "url": "https://www.georisques.gouv.fr",
    "attribution": "Source Georisques (BRGM / MTE) - donnees a jour au <millesime>",
    "recuperee_le": "2026-08-25T10:00:00+00:00"
  },
  "erreurs_partielles": [],
  "genere_le": "2026-08-25T10:00:01+00:00",
  "avertissement": "Ce rapport agrege les donnees publiques Georisques..."
}
\end{lstlisting}

\section{Champs interdits}

Les champs suivants ne peuvent pas exister dans le contrat. Leur présence fait
échouer la validation du schéma, et une porte de recherche automatique échoue si
l'un d'eux apparaît dans un schéma émis :

\begin{itemize}[leftmargin=1.4em]
  \item \code{niveau} et les bandes D03 ;
  \item tout score numérique, agrégé ou composite ;
  \item \code{recommandations} et toute évaluation générée par IA ;
  \item \code{copernicus} dans le contrat canonique v1.
\end{itemize}

\section{Règles de sérialisation}

\begin{table}[H]
\centering
\small
\caption{Règles de sérialisation à respecter par tout consommateur.}
\label{tab:serialisation}
\begin{tabularx}{\textwidth}{@{}l X@{}}
\toprule
\textbf{Règle} & \textbf{Énoncé} \\
\midrule
Optionnels toujours présents &
  Les champs spécifiques à un aléa sont toujours présents, à \code{null}
  lorsqu'ils ne s'appliquent pas. Une clé n'est jamais absente. \\
\addlinespace
Exception \code{per\_building} &
  L'objet \code{per\_building} n'existe que si le point a été testé
  géométriquement. Sa présence porte l'information. \\
\addlinespace
BDNB verbatim &
  Les clés BDNB reçues et émises sont identiques — vérifié par test de
  conformité. \\
\addlinespace
Attribution non vide &
  Chaque bloc \code{\_source} porte une chaîne d'attribution non vide. \\
\addlinespace
Résolution contrainte &
  Le RGA ne peut sérialiser aucune autre valeur que
  \code{commune-level-estimate}. \\
\bottomrule
\end{tabularx}
\end{table}

\section{Le référentiel des treize aléas}

\begin{table}[H]
\centering
\footnotesize
\caption{Les treize aléas, leur libellé et leur niveau de résolution attendu.}
\label{tab:aleas}
\begin{tabularx}{\textwidth}{@{}l l X@{}}
\toprule
\textbf{Code} & \textbf{Libellé} & \textbf{Résolution attendue} \\
\midrule
\code{icpe} & Installations classées &
  \code{per-building} ou \code{commune-level} \\
\code{inondation} & Inondation & \code{per-building} \\
\code{sismicite} & Sismicité & \code{commune-level} (décret, permanent) \\
\code{mouvement\_terrain} & Mouvement de terrain & \code{per-building} \\
\code{radon} & Radon & \code{commune-level} (permanent) \\
\code{rga} & Retrait-gonflement des argiles &
  \code{commune-level-estimate} (forcé) \\
\code{cavite} & Cavités souterraines & \code{per-building} ou communal \\
\code{feu\_foret} & Feu de forêt & \code{per-building} \\
\code{avalanche} & Avalanche & \code{per-building} \\
\code{canalisations} & Canalisations (TMD) &
  \code{per-building} par proximité (200~m) \\
\code{vent\_cyclonique} & Vent cyclonique & \code{commune-level} \\
\code{ppr} & PPR (agrégat) & \code{per-building} \\
\code{ssp} & Sites et sols pollués &
  \code{per-building} par proximité (200~m) \\
\bottomrule
\end{tabularx}
\end{table}
